\chapter{CNN nhận dạng hình học}
\label{chuong:hinhhoc}

\section{Dữ liệu}
Không có bộ dữ liệu công khai phù hợp nên dữ liệu hình học được \textbf{sinh bằng chương trình}: bộ rasterizer giải tích bằng NumPy với siêu lấy mẫu $4\times4$, mỗi ảnh gốc $28\times28$ chứa một hình có
độ dày nét 1--3 px, góc quay $0^{\circ}$--$360^{\circ}$, đường kính ngoại tiếp 7--21 px, nhiễu Gauss nhẹ và nhiễu muối tiêu. Ba lớp: \emph{tròn}, \emph{tam giác}, \emph{chữ nhật hoặc vuông}.

Để mạng không chỉ học các hình quá lý tưởng, bộ sinh được làm đa dạng:
\begin{table}[H]\centering
\begin{tabular}{lp{0.66\textwidth}}
\toprule
Lớp & Các dạng được sinh \\
\midrule
Tam giác & đều 25\%, vuông 20\%, cân 20\%, bất kỳ 35\%; mọi góc trong $\ge18^{\circ}$ \\
Chữ nhật & tỉ lệ cạnh $U(0{,}3;1)$, 25\% là hình vuông; cạnh ngắn $\ge$ độ dày nét $+2$ px để lõi không bị lấp đặc \\
Tròn & 50\% elip (tỉ lệ trục $U(0{,}7;1)$), 50\% lượn sóng bán kính tới 5\%, phần còn lại tròn hoàn hảo \\
Rung tay & 50\% tam giác/chữ nhật có nhiễu Gauss lên đỉnh, $\sigma$ tới 2,5\% kích thước \\
\bottomrule
\end{tabular}
\caption{Độ đa dạng của bộ sinh dữ liệu hình học.}
\end{table}

Tổng cộng 13\,500 ảnh chia phân tầng 70/15/15 thành 9\,450 train, 2\,025 validation, 2\,025 test (mỗi lớp 675 ảnh test), các tập không giao nhau và không có ảnh trùng.
Ngoài ra có bộ \emph{stress} 300 ảnh khó hơn (nhiễu mạnh, mờ, có cột bị mất, đỉnh rung 3{,}5\%, hình tròn luôn méo) chỉ dùng để đo độ bền, không dùng huấn luyện hay chọn mô hình.
Mỗi ảnh gốc được đưa qua \emph{chính} tiền xử lý của hệ thống (Chương~\ref{chuong:tienxuly}) để tạo đầu vào CNN, nên huấn luyện và nhận dạng dùng chung một đường xử lý.

\hinh{hinh_hoc_xem_truoc.png}{0.75}{Ảnh hình học tổng hợp (gốc $28\times28$).}

\section{Kiến trúc và huấn luyện}
Cùng bộ khung với mạng chữ số nhưng đầu ra 3 lớp. Cấu hình \emph{cơ sở}: không BatchNorm, không Dropout, không tăng cường dữ liệu (56\,259 tham số). Cấu hình \emph{cải tiến}: thêm BatchNorm \cite{ioffe2015} sau hai
tầng tích chập, Dropout $0{,}3$ \cite{srivastava2014} và tăng cường dữ liệu (xoay $\pm20^{\circ}$, dịch $\pm1{,}5$ px, co giãn, đổi độ sáng, nhiễu Gauss $0{,}012$) (56\,355 tham số).
Adam $10^{-3}$, batch 64, 20 epoch, hạt 42, CPU; chọn epoch và cấu hình cuối bằng validation như ở Chương~\ref{chuong:chuso}.
Cấu hình cải tiến thay đồng thời ba yếu tố (BatchNorm, Dropout, tăng cường) nên không quy được khác biệt cho riêng yếu tố nào.

\section{Kết quả}
\begin{table}[H]\centering
\begin{tabular}{lrrrr}
\toprule
Cấu hình & Validation & Test (2\,025 ảnh) & F1 macro & Stress (300 ảnh) \\
\midrule
Cơ sở & $99{,}75\%$ & $2014/2025=99{,}46\%$ & $99{,}46\%$ & $258/300=86{,}0\%$ \\
Cải tiến (được chọn) & $99{,}90\%$ & $2019/2025=99{,}70\%$ & $99{,}70\%$ & $270/300=90{,}0\%$ \\
\bottomrule
\end{tabular}
\caption{Kết quả hình học; cấu hình cuối chọn bằng validation. Test và stress là dữ liệu tổng hợp.}
\end{table}
Khoảng Wilson 95\% của accuracy test: cơ sở $[99{,}03;\,99{,}70]\%$, cải tiến $[99{,}36;\,99{,}86]\%$. McNemar chính xác ghép cặp: cơ sở sai/cải tiến đúng 6 ảnh, ngược lại 1 ảnh, $p=0{,}125$
(chưa đủ để khẳng định khác biệt). Theo lớp, cấu hình được chọn có precision/recall từ $99{,}56\%$ đến $100\%$.
Sáu ảnh test sai gồm hai elip, một elip lượn sóng, hai chữ nhật và một hình vuông; các nhầm lẫn là tròn$\to$chữ nhật (3) và chữ nhật$\to$tam giác (3).

\hinh{hinh_hoc_ma_tran.png}{0.5}{Ma trận nhầm lẫn của cấu hình được chọn trên tập test hình học.}

Trên bộ stress, các nhầm lẫn lớn nhất là chữ nhật$\to$tam giác (13 ảnh) và tròn$\to$tam giác (9), tròn$\to$chữ nhật (7). Xem lại 24 lỗi có softmax cao nhất cho thấy phần lớn là hình \emph{bị đứt nét} (bộ stress xóa một cột pixel): bước giữ thành phần liên thông lớn nhất chỉ còn một cung hoặc một góc, và mạng đoán sai trên phần còn lại. Đây là hệ quả của giới hạn của bước liên thông chứ không chỉ của mạng.

\section{Độ bền theo góc xoay và độ dày nét}
Quét ghép cặp 40 đối tượng gốc $\times$ 12 góc $\times$ 3 độ dày nét (tam giác và chữ nhật của bộ stress, 480 ảnh mỗi độ dày, mỗi mô hình), thực hiện sau khi đã chọn mô hình:
\begin{table}[H]\centering
\begin{tabular}{lrrr}
\toprule
Độ dày nét (px) & 1 & 2 & 3 \\
\midrule
Cơ sở & $90{,}8\%$ & $92{,}1\%$ & $89{,}0\%$ \\
Cải tiến & $94{,}2\%$ & $94{,}8\%$ & $95{,}0\%$ \\
\bottomrule
\end{tabular}
\caption{Độ chính xác trong quét ghép cặp (dữ liệu tổng hợp, không phải hình người vẽ).}
\end{table}

\hinh{hinh_hoc_quet.png}{0.95}{Độ chính xác theo góc xoay và độ dày nét (mỗi ô 40 đối tượng khớp nhau).}

\section{So sánh dữ liệu huấn luyện}
\label{muc:sosanhdulieu}
Để kiểm tra việc làm đa dạng bộ sinh có thật sự cần thiết, cấu hình cải tiến được huấn luyện thêm trên hai bộ dữ liệu khác với cùng kiến trúc, công thức, hạt và tiền xử lý:
\emph{đơn giản} (tam giác đều, chữ nhật tỉ lệ 0,55--1, hình tròn hoàn hảo, không rung đỉnh; 4\,500 ảnh) và \emph{đa dạng} với cùng 4\,500 ảnh. Cả ba mô hình (kèm mô hình chính 13\,500 ảnh) được đo trên các bộ giống nhau.

\begin{table}[H]\centering\small
\begin{tabular}{lrrrr}
\toprule
Dữ liệu huấn luyện & Test đơn giản & Test đa dạng & Stress đa dạng & Hình vẽ Pillow (TB) \\
\midrule
Đơn giản, 4.500 ảnh & 100,0\% & 96,9\% & 80,3\% & 91,5\% \\
Đa dạng, 4.500 ảnh & 99,9\% & 99,5\% & 85,3\% & 99,7\% \\
Đa dạng, 13.500 ảnh (mô hình chính) & 100,0\% & 99,7\% & 90,0\% & 99,6\% \\
\bottomrule
\end{tabular}
\caption{Ảnh hưởng của dữ liệu huấn luyện (tất cả là dữ liệu tổng hợp).}
\end{table}

\begin{table}[H]\centering\small
\begin{tabular}{lrrr}
\toprule
Họ hình (Pillow) & Đơn giản 4\,500 & Đa dạng 4\,500 & Đa dạng 13\,500 \\
\midrule
Tròn & 100,0\% & 100,0\% & 100,0\% \\
Tam giác đều & 100,0\% & 100,0\% & 100,0\% \\
Tam giác vuông & 100,0\% & 100,0\% & 99,4\% \\
Tam giác bất kỳ & 100,0\% & 98,1\% & 99,4\% \\
Chữ nhật (0,55--1) & 100,0\% & 100,0\% & 100,0\% \\
Chữ nhật dài (0,25--0,45) & 48,8\% & 100,0\% & 98,8\% \\
\bottomrule
\end{tabular}
\caption{Độ chính xác theo họ hình trên ảnh vẽ Pillow (trung bình 4 độ dày nét, 160 ảnh mỗi họ).}
\end{table}
\hinh{so_sanh_du_lieu_hinh.png}{0.85}{So sánh dữ liệu huấn luyện hình học.}

Bộ test của chính dữ liệu đơn giản đã bão hòa ($100\%$) nên không phân biệt được mô hình tổng quát với mô hình học thuộc phân bố hẹp; trên test đa dạng mô hình đó chỉ đạt $96{,}9\%$.
Dữ liệu đa dạng cải thiện rõ nhất ở \textbf{chữ nhật dài} ($48{,}8\%\to100\%$ với 4\,500 ảnh, $98{,}8\%$ với 13\,500 ảnh); các họ còn lại gần như không đổi, vì tam giác vuông và tam giác bất kỳ
đã được nhận đúng ngay cả khi mô hình chỉ học tam giác đều (với tiền xử lý hiện tại). Gấp ba số ảnh tăng độ chính xác trên bộ stress ($85{,}3\%\to90{,}0\%$) nhưng không thấy khác biệt trên hình vẽ Pillow ($99{,}7\%$ so với $99{,}6\%$, trong sai số của 960 ảnh).
Một hạt cho mỗi cấu hình, nên các chênh lệch nhỏ chưa đủ ý nghĩa thống kê.

\section{Hạn chế}
Mọi số liệu ở chương này là dữ liệu tổng hợp cùng họ bộ sinh; bộ stress và bộ quét cũng do cùng bộ sinh tạo ra nên accuracy cao \textbf{không bảo đảm} hiệu quả trên hình người vẽ.
Hình tròn méo nặng kèm nhiễu mạnh và chữ nhật quá dài vẫn dễ sai: trong phép thử chữ nhật nét 3--8 px, tỉ lệ cạnh 0,22 và 0,31 nhận đúng, từ 0,15 trở xuống (ngoài dải huấn luyện 0,3--1) bị đoán thành tam giác. Mạng chỉ có ba lớp và không từ chối ngôi sao, hình thang hay chữ cái.
